\documentclass[11pt,a4paper]{article}
\usepackage[utf8]{inputenc}\usepackage[T1]{fontenc}
\usepackage[provide=*,english]{babel}
\usepackage{lmodern,microtype,amsmath,amssymb,booktabs,array,geometry,xcolor,fancyhdr}
\usepackage[hidelinks]{hyperref}
\geometry{margin=2.3cm,headheight=14pt}
\setlength{\parindent}{0pt}\setlength{\parskip}{0.4em}
\definecolor{gcblue}{RGB}{34,73,110}\definecolor{gcgray}{RGB}{90,96,102}
\pagestyle{fancy}\fancyhf{}
\fancyhead[L]{\small\color{gcgray}GC2D: BM4 family}
\fancyhead[R]{\small\color{gcgray}Model theory}\fancyfoot[C]{\thepage}
\newcommand{\R}{\mathbb{R}}
\begin{document}
\begin{center}
 {\LARGE\bfseries\color{gcblue}BM4 Fourth-Order Composition}\par
 \vspace{0.35em}{\large Direct--adjoint splitting and all projection configurations}
\end{center}

\section{Common composition}

Let $\chi_h$ be a first-order map and $\chi_h^*=\chi_{-h}^{-1}$ its adjoint.
All BM4 configurations use the palindromic twelve-stage composition
\[
\Psi_h=\chi_{b_{12}h}\circ\chi^*_{b_{11}h}\circ\cdots
\circ\chi_{b_2h}\circ\chi^*_{b_1h},
\]
with alternating adjoint and direct maps in chronological execution order and
\[
(b_1,\ldots,b_{12})=(a_1,a_2,a_3,a_4,a_5,a_6,a_6,a_5,a_4,a_3,a_2,a_1),
\]
\begin{align*}
a_1&=0.0792036964311957,&a_2&=0.1303114101821663,\\
a_3&=0.2228614958676077,&a_4&=-0.3667132690474257,\\
a_5&=0.3246481886897062,&a_6&=0.1096884778767498.
\end{align*}
The half-cycle coefficients sum to $1/2$.  Palindromy makes the composition
self-adjoint and its order conditions give designed global order four when the
direct--adjoint pair satisfies the assumed smoothness and consistency.  The
negative $a_4$ is an intentional backward subflow.

\section{Public configurations}

The family exposes six classes because projection placement changes the
numerical map:
\begin{center}
\small
\begin{tabular}{@{}p{0.27\linewidth}p{0.48\linewidth}p{0.17\linewidth}@{}}
\toprule
Class & projection placement & nonlinear unknown\\
\midrule
\texttt{BM4Composition} & formulation-owned output extraction & none\\
\texttt{ProjectedBM4}\allowbreak\texttt{Composition} & arithmetic re-embedding after every stage & none\\
\texttt{MidpointBM4} & arithmetic re-embedding after the full cycle & none\\
\texttt{BM4Implicit1} & one symmetric projection around the full cycle & $\mu\in\R^{2N}$\\
\texttt{BM4Implicit2} & the same projection in simultaneous form & $(Y,\mu)\in\R^{6N}$\\
\texttt{BM4\_implicit2} & full-diagonal projection of duplicated $Z$ & $\mu\in\R^4$\\
\bottomrule
\end{tabular}
\end{center}
The underscore in \texttt{BM4\_implicit2} distinguishes the fully extended
method from the simultaneous physical method \texttt{BM4Implicit2}.

\section{Doubled physical state}

For a physical state $z\in\R^m$, define
\[
E=\begin{pmatrix}I\\I\end{pmatrix},\quad
P=\tfrac12\begin{pmatrix}I&I\end{pmatrix},\quad
G=\begin{pmatrix}I&-I\end{pmatrix},\quad N=G^T.
\]
The unprojected composition retains two copies.  Stage projection applies
$EP$ after each of the twelve factors; midpoint projection applies it once
after the complete cycle.  These placements must not be conflated because
projection and the stage maps do not commute.

\section{Implicit physical projection}

Let $\Psi_{h,t}$ denote the complete coupled BM4 cycle.  The reduced method
solves
\begin{align}
M(\mu)&=\Psi_{h,t}(Ez+N\mu),\\
r(\mu)&=GM(\mu)+2\mu=0,\\
z^+&=P\bigl(M(\mu)+N\mu\bigr).
\end{align}
With $J_{\rm BM4}=D\Psi_{h,t}(Ez+N\mu)$, the exact Newton matrix is
\[
D_\mu r=GJ_{\rm BM4}N+2I.
\]
The simultaneous method solves
\[
F(Y,\mu)=
\begin{pmatrix}
Y-N\mu-\Psi_{h,t}(Ez+N\mu)\\GY
\end{pmatrix}=0,
\]
whose block Jacobian is
\[
\begin{pmatrix}I&-(I+J_{\rm BM4})N\\G&0\end{pmatrix}.
\]
Its Schur complement is the reduced Newton equation, proving equivalence of
the two ideal roots.  Both formulations accept exact Newton or good Broyden;
this solver choice changes work and finite-tolerance behavior, not the
defining projected map.

\section{Fully extended configuration}

For one non-autonomous guiding-centre particle introduce
\[
Z=(x,y,t,k),\qquad K(Z)=H(t,x,y)+k.
\]
\texttt{BM4\_implicit2} duplicates $Z$ to an eight-dimensional splitting
state, executes the same twelve-stage coefficient schedule, and projects the
full diagonal with a four-component multiplier.  The accepted physical
solution contains $(x,y)$ while diagnostics retain $t$, $k$, and the drift of
$K$.  This configuration is a genuine extended-state map and is not algebraic
workspace for \texttt{BM4Implicit2}.

\section{Geometric interpretation and finite solves}

For an exact adjoint pair, the base composition is symmetric and fourth order.
The symmetric implicit projection is designed to preserve the inherited
geometric structure at its exact root.  Arithmetic stage or cycle projections
define separate maps and must be assessed directly.  Finite Newton or Broyden
tolerances perturb every ideal statement; the implementation therefore records
nonlinear work and exposes step maps for Jacobian-based audits.

\section{Implementation correspondence}

The shared composition is in the BM4 module's \texttt{\_core.py}; the midpoint,
reduced, simultaneous, and fully extended variants are in its neighboring
modules.  Two detailed derivations accompany this theory:
\begin{itemize}
\item \texttt{implicit-reduced.tex}; and
\item \texttt{implicit-simultaneous.tex}.
\end{itemize}
The runtime architecture lives under \texttt{../simulation/}; shared physical
dynamics are documented outside the model tree.

\end{document}
